\documentclass{article}
\usepackage{iclr2027_conference,times}
\renewcommand{\iclrruler}[1]{}
\renewcommand{\headrulewidth}{0pt}
\input{math_commands.tex}
\usepackage{amsmath,amssymb,amsthm,mathtools}
\usepackage{booktabs}
\usepackage{multirow}
\usepackage{graphicx}
\usepackage{microtype}
\usepackage{xspace}
\usepackage{hyperref}
\usepackage{url}

\newtheorem{theorem}{Theorem}
\newtheorem{lemma}{Lemma}
\newtheorem{proposition}{Proposition}
\newtheorem{corollary}{Corollary}
\newtheorem{definition}{Definition}
\newcommand{\method}{\textsc{LogRecal}\xspace}
\newcommand{\calK}{\mathcal{K}}
\newcommand{\Ber}{\operatorname{Ber}}
\newcommand{\simplex}{\Delta}

\title{Near-Optimal Online Recalibration for Log Loss}
\author{Pahan Dewasurendra \\ Johns Hopkins University \\ July 23, 2026}

% \iclrfinalcopy
\begin{document}
\maketitle

\begin{abstract}
Online recalibration replaces an arbitrary forecast sequence by predictions
that are calibrated and nearly as accurate. Recent work attains the optimal
$T=\Theta(\varepsilon^{-3})$ horizon for globally Lipschitz proper losses,
but logarithmic loss is unbounded at probabilities zero and one. Naive
clipping makes the existing bound polynomially worse. We give the first
near-optimal recalibration guarantee for Bernoulli log loss without any
restriction on the input forecasts. Our algorithm achieves expected ordinary
$\ell_1$ calibration at most $\varepsilon$ and expected excess log loss at
most $\varepsilon^2$ after
$O(\varepsilon^{-3}\log(1/\varepsilon))$ rounds. A matching
$\Omega(\varepsilon^{-3})$ lower bound leaves only one logarithmic factor.
The key is a nonuniform mixture oracle controlled by adjacent KL tangent
gaps, rather than by a global Lipschitz constant. A grid uniform in the
Bernoulli information coordinate $\arcsin\sqrt p$ has
$O(\varepsilon^{-1})$ points, probability mesh $O(\varepsilon)$, and
labelwise tangent gap $O(\varepsilon^2)$ even at the boundary. Normalizing
the loss objective and initializing the imbalanced meta learner at that
objective prevents the growing clipped loss range from being paid twice.
Experiments on controlled streams and a shifted CIFAR-10 binary stream
evaluate the finite-time calibration and loss frontier. A fixed practical
step choice reduces exact calibration error by $92.8\%$ on the shifted
image stream while improving log loss. The result closes
the unbounded-log-loss gap for arbitrary-hint online recalibration while
making explicit which information-geometric ingredients are inherited from
earlier log-loss discretization.
\end{abstract}

\section{Introduction}

Probability forecasts are often optimized for a proper loss, then adjusted
to improve calibration. In the online formulation, an arbitrary hint
$q_t\in[0,1]$ arrives on round $t$. A recalibrator replaces it by
$p_t\in[0,1]$ before observing the binary outcome $y_t$. Its two goals are
different. Predictions with the same reported probability should have the
right empirical frequency, while their average loss should remain close to
that of the hints.

Kuleshov and Ermon introduced this adversarial recalibration problem
\citep{kuleshov2017estimating}. Blackwell-style methods later improved its
finite-time rates \citep{okoroafor2024faster}. Most recently,
\citet{hu2026optimal} proved that $(\varepsilon,\varepsilon^2)$
recalibration is possible in $\Theta(\varepsilon^{-3})$ rounds for
globally Lipschitz proper losses. This is the statistically natural target:
calibration is controlled at scale $\varepsilon$, while the loss cost of a
local probability correction is controlled at scale $\varepsilon^2$.

Log loss is the central exception. For $y\in\{0,1\}$,
\begin{equation}
  \ell_{\log}(p,y)=-y\log p-(1-y)\log(1-p).
  \label{eq:logloss}
\end{equation}
Its derivative and range diverge at the probability boundary.
\citet{hu2026optimal} explicitly leave open whether their Lipschitz
assumption is necessary and mention clipping as a practical approximation.
To tolerate arbitrary hints, clipping must reach
$\delta=O(\varepsilon^2)$. The resulting Lipschitz constant is
$\Theta(\varepsilon^{-2})$, so direct use of their theorem gives the
uninformative horizon $O(\varepsilon^{-7})$.

The obstacle is not the total range of log loss. It is the local geometric
error paid by the oracle that balances calibration against loss. A
probability-uniform grid has badly distorted cells near zero and one. In the
coordinate
\begin{equation}
  \theta(p)=\arcsin\sqrt p,
  \label{eq:fisher-coordinate}
\end{equation}
Bernoulli Fisher information is constant up to scale. Equal angular cells
therefore shrink toward the boundary. This grid is classical
\citep{derooij2009switching,kotlowski2016isotonic} and has recently been
used for KL calibration \citep{luo2025kl,liu2026adaptive}. Our contribution
is not the grid. It is a two-objective oracle and analysis that make this
geometry compatible with arbitrary hints and ordinary calibration.

The oracle mixes at most two adjacent grid points. Its calibration error is
the probability width of the cell. Its loss error is an oriented tangent
gap, which is a normalized adjacent Bernoulli KL divergence. On the grid in
Equation~(\ref{eq:fisher-coordinate}), these two errors are respectively
$O(\varepsilon)$ and $O(\varepsilon^2)$ using only
$O(\varepsilon^{-1})$ actions. Clipping then costs only
$O(\varepsilon^2)$ in the one-sided comparison with the original hint.

One further issue remains. The clipped log-loss range is
$M=O(\log(1/\varepsilon))$. A generic range-based approachability bound
would pay $M^2$. We scale the loss coordinate by $M$ and initialize the
two-objective meta learner at the loss vertex. The loss comparator then has
zero initial divergence. Choosing its step size at
$\Theta(\varepsilon^2/M)$ pays $M$ only once.

Our contributions are:
\begin{enumerate}
  \item We give a nonuniform mixture linear optimization oracle whose
  calibration tolerance is the grid mesh and whose proper-loss tolerance is
  an adjacent tangent gap. This replaces global Lipschitz control by local
  Bayes regret.
  \item We instantiate the oracle for log loss with a Fisher-uniform grid.
  For arbitrary hints, including zero and one, the resulting algorithm
  achieves expected $(\varepsilon,\varepsilon^2)$ recalibration in
  $O(\varepsilon^{-3}\log(1/\varepsilon))$ rounds.
  \item We transfer the recent recalibration hard instance through Pinsker's
  inequality to prove an $\Omega(\varepsilon^{-3})$ lower bound for log
  loss.
  \item We test the algorithm on boundary-sensitive controlled sequences and
  on a public binary image-classification stream under distribution shift.
  The evaluations report ordinary finite-grid calibration and realized log
  excess, rather than only the mixed-action surrogate used by the oracle.
\end{enumerate}

\section{Problem and background}
\label{sec:problem}

\paragraph{Protocol.}
At round $t$, an environment reveals a hint $q_t\in[0,1]$. The learner then
randomly predicts $p_t\in[0,1]$, after which $y_t\in\{0,1\}$ is revealed.
The environment may depend on the past, but it may not observe the learner's
fresh random draw before selecting $y_t$. Expectations below are only over
learner randomization unless a stochastic lower-bound construction is stated.

For a finite-valued prediction sequence, define ordinary $\ell_1$
calibration as
\begin{equation}
  \calK_1(p_{1:T},y_{1:T})
  =\frac1T\sum_{s\in\operatorname{range}(p_{1:T})}
  \left|\sum_{t=1}^T \mathbf{1}\{p_t=s\}(s-y_t)\right|.
  \label{eq:k1}
\end{equation}
This definition charges each reported probability separately. In
particular, it is not a binned diagnostic and does not merge nearby values.

\begin{definition}[Log-loss recalibration]
For $\varepsilon\in(0,1/4)$, a prediction sequence is
$(\varepsilon,\varepsilon^2)$-recalibrated against
$(q_{1:T},y_{1:T})$ if
\begin{align}
  \mathbb{E}\calK_1(p_{1:T},y_{1:T})&\leq\varepsilon,
  \label{eq:target-cal}\\
  \mathbb{E}\frac1T\sum_{t=1}^T
  [\ell_{\log}(p_t,y_t)-\ell_{\log}(q_t,y_t)]
  &\leq\varepsilon^2.
  \label{eq:target-loss}
\end{align}
The second inequality is one-sided. At boundary hints, it is interpreted in
the extended-real sense.
\end{definition}

\paragraph{Simultaneous approachability.}
We build on the imbalanced simultaneous Blackwell framework of
\citet{hu2026optimal}, itself based on simultaneous approachability
\citep{hu2026simultaneous} and classical Blackwell approachability
\citep{blackwell1956analog}. For a finite grid $\mathcal G$, an action
$a\in\simplex^{\mathcal G}$ is a distribution over reports. The two payoff
coordinates are
\begin{align}
 v^{(1)}(a,y)&=\{a_s(s-y)\}_{s\in\mathcal G},
 \label{eq:cal-payoff}\\
 v^{(2)}(a,q,y)&=\frac1M\sum_{s\in\mathcal G}a_s
 [\ell_{\log}(s,y)-\ell_{\log}(\bar q,y)],
 \label{eq:loss-payoff}
\end{align}
where $\bar q$ is a clipped hint and $M$ is the clipped loss range.
The first coordinate is tested by
$u\in[-1,1]^{\mathcal G}$. Its support function is exactly
Equation~(\ref{eq:k1}) when the realized reports are point masses. The second
coordinate has the singleton test $1$.

Given weights $(\alpha,\beta)\in\simplex^2$, a mixture linear optimization
oracle must choose one action that approximately satisfies both half-space
constraints:
\begin{equation}
 \alpha\langle u,v^{(1)}(a,y)\rangle
 +\beta v^{(2)}(a,q,y)
 \leq\alpha\epsilon_1+\beta\epsilon_2
 \quad\text{for both }y\in\{0,1\}.
 \label{eq:mloo-goal}
\end{equation}
The different tolerances are essential. We need
$\epsilon_1=O(\varepsilon)$ and
$M\epsilon_2=O(\varepsilon^2)$.

\section{A curvature-adaptive mixture oracle}
\label{sec:oracle}

Let
$\mathcal G=\{p_0<\cdots<p_N\}\subset(0,1)$ be symmetric about $1/2$.
Define its extended probability mesh
\begin{equation}
 h(\mathcal G)=\max\{p_0,1-p_N,\max_{i<N}(p_{i+1}-p_i)\}.
 \label{eq:mesh}
\end{equation}
For adjacent $p<r$, define the oriented log tangent gap
\begin{equation}
 \tau(p,r)=
 \begin{cases}
 \KL(\Ber(r)\Vert\Ber(p))/(1-r),&r\leq1/2,\\
 \KL(\Ber(p)\Vert\Ber(r))/p,&p\geq1/2,\\
 \max\{\KL(\Ber(p)\Vert\Ber(r))/p,
 \KL(\Ber(r)\Vert\Ber(p))/(1-r)\},&p<1/2<r.
 \end{cases}
 \label{eq:tangent-gap}
\end{equation}
Set $\tau(\mathcal G)=\max_{i<N}\tau(p_i,p_{i+1})$.
The orientation in Equation~(\ref{eq:tangent-gap}) is selected by the side
of $1/2$ containing the cell. It comes from the slope of the proper-loss
support line, not from a symmetric smoothness bound.

\begin{lemma}[Nonuniform log-loss oracle]
\label{lem:oracle}
Clip $q$ to $\bar q\in[p_0,p_N]$ and let
$M=\log(1/p_0)$ for a symmetric grid. For every
$u\in[-1,1]^{\mathcal G}$ and
$(\alpha,\beta)\in\simplex^2$, there is a distribution $a$ supported on one
point or two adjacent points such that, for both labels,
\begin{equation}
 \alpha\langle u,v^{(1)}(a,y)\rangle+
 \beta v^{(2)}(a,\bar q,y)
 \leq \alpha h(\mathcal G)+
 \beta\frac{\tau(\mathcal G)}M.
 \label{eq:oracle-bound}
\end{equation}
\end{lemma}

The proof modifies the adjacent-line construction of
\citet{okoroafor2024faster,hu2026optimal}. For $\alpha>0$, define
\begin{equation}
 d(p,y)=\frac{\beta}{\alpha M}
 [\ell_{\log}(p,y)-\ell_{\log}(\bar q,y)]
 \quad\text{and}\quad
 f_i=(u_i p_i+d(p_i,0),u_i(p_i-1)+d(p_i,1)).
 \label{eq:f-curve}
\end{equation}
If a point is already below the target quadrant, a point mass works.
Otherwise, the sequence contains an adjacent upper-left to lower-right
crossing. For a crossing $p<r$ on the upper half, the vertical gap from the
log-loss curve at $r$ to the support line at $p$ is exactly
\begin{equation}
 \frac{\beta}{\alpha M}
 \left[
 \log\frac p r+\frac{1-p}{p}\log\frac{1-p}{1-r}
 \right]
 =
 \frac{\beta}{\alpha M}
 \frac{\KL(\Ber(p)\Vert\Ber(r))}{p}.
 \label{eq:gap-identity}
\end{equation}
Reflection gives the first case of Equation~(\ref{eq:tangent-gap}).
Changing the calibration direction between adjacent points costs at most
$r-p$. The two support lines then determine a mixture whose two label
payoffs obey Equation~(\ref{eq:oracle-bound}). The endpoint residuals are at
most $p_0$ and $1-p_N$, which explains the extended mesh. The case
$\alpha=0$ follows by compactness. Appendix~\ref{app:oracle-proof} gives the
full argument and a general proper-loss identity.

\paragraph{Beyond log loss.}
The oracle is not specific to KL divergence. If a differentiable binary
proper loss has Bayes-risk weight $w(p)=-H''(p)$, replace the two oriented
gaps in Equation~(\ref{eq:tangent-gap}) by
\begin{equation}
 \frac1p\int_p^r(t-p)w(t)\,dt
 \quad\text{and}\quad
 \frac1{1-r}\int_p^r(r-t)w(t)\,dt.
 \label{eq:general-gaps-main}
\end{equation}
The same mixture oracle holds. This gives a grid-dependent sufficient
condition through local Bayes regret, with no global Lipschitz requirement.
Log loss is the important unbounded case where one grid simultaneously
controls these gaps, ordinary calibration resolution, clipping cost, and
action count.

\subsection{The Fisher grid}

Fix an angular resolution $a\in(0,\pi/8]$. Choose
\begin{equation}
 N=\left\lceil\frac{\pi/2-2a}{a}\right\rceil,\qquad
 \theta_i=a+i\frac{\pi/2-2a}{N},\qquad
 p_i=\sin^2\theta_i,
 \label{eq:grid}
\end{equation}
with equal angular spacing and symmetric endpoints.
This construction is uniform in half the usual Fisher arc-length
coordinate. It was introduced for Bernoulli log-loss discretization by
\citet{derooij2009switching} and reused in online isotonic regression and KL
calibration \citep{kotlowski2016isotonic,luo2025kl,liu2026adaptive}.

\begin{lemma}[Grid geometry]
\label{lem:grid}
The grid in Equation~(\ref{eq:grid}) satisfies
\begin{equation}
 |\mathcal G|=O(a^{-1}),\qquad
 h(\mathcal G)\leq a,\qquad
 p_0=\sin^2a\leq a^2,\qquad
 \tau(\mathcal G)\leq72a^2.
 \label{eq:grid-bounds}
\end{equation}
\end{lemma}

The mesh bound follows because
$\frac{d}{d\theta}\sin^2\theta=\sin2\theta\leq1$.
For the tangent gap, let
$p=\sin^2\theta<r=\sin^2\phi\leq1/2$ and
$s=\phi-\theta\leq a\leq\theta$. The Bernoulli chi-squared bound gives
\begin{align}
 \frac{\KL(\Ber(r)\Vert\Ber(p))}{1-r}
 &\leq\frac{(r-p)^2}{p(1-p)(1-r)} \nonumber\\
 &=\frac{\sin^2(\phi+\theta)\sin^2s}
 {\sin^2\theta\cos^2\theta\cos^2\phi}
 \leq72s^2.
 \label{eq:grid-proof}
\end{align}
The upper half follows by reflection. A central cell has both Bernoulli
variances bounded below and satisfies the same loose constant.

\paragraph{Clipping is only second order.}
Let $\delta=p_0=\sin^2a$ and
$\bar q=\min\{\max\{q,\delta\},1-\delta\}$. For either label,
\begin{equation}
 \ell_{\log}(\bar q,y)-\ell_{\log}(q,y)
 \leq-\log(1-\delta)=O(a^2).
 \label{eq:clip}
\end{equation}
The bound is one-sided, exactly as required for excess loss. No lower bound
on the original hints is assumed.

\section{Near-optimal log-loss recalibration}
\label{sec:algorithm}

\method instantiates imbalanced simultaneous approachability with
Lemma~\ref{lem:oracle}. It maintains a calibration distinguisher
$u_t\in[-1,1]^{\mathcal G}$ and meta weights
$w_t=(\alpha_t,\beta_t)\in\simplex^2$. The initialization is
\begin{equation}
 u_1=0,\qquad w_1=(0,1).
 \label{eq:init}
\end{equation}
On each round, the oracle returns a sparse distribution $a_t$. The learner
samples $p_t\sim a_t$ with fresh randomness before observing $y_t$.
Afterward it performs projected gradient ascent on the calibration test:
\begin{equation}
 u_{t+1}=\Pi_{[-1,1]^{\mathcal G}}
 \left(u_t+\zeta\{\mathbf{1}\{p_t=s\}(s-y_t)\}_{s\in\mathcal G}\right).
 \label{eq:u-update}
\end{equation}
The meta learner receives the two constraint residuals and performs
Euclidean ascent on $\simplex^2$:
\begin{equation}
 w_{t+1}=\Pi_{\simplex^2}(w_t+\gamma g_t).
 \label{eq:meta-update}
\end{equation}
For reduced variance, $g_t$ can use the known mixed payoff under $a_t$.
Its conditional expectation equals the point-mass payoff required by the
approachability proof. Algorithm~\ref{alg:logrecal} is written explicitly in
Appendix~\ref{app:algorithm}.

\begin{theorem}[Expected log-loss recalibration]
\label{thm:upper}
There are universal constants $C,c>0$ such that the following holds.
For any $\varepsilon\in(0,c)$, any horizon
\begin{equation}
 T\geq C\frac{\log(1/\varepsilon)}{\varepsilon^3},
 \label{eq:horizon}
\end{equation}
and every non-anticipating hint-label sequence, \method produces predictions
satisfying Equations~(\ref{eq:target-cal}) and
(\ref{eq:target-loss}).
\end{theorem}

\paragraph{Proof overview.}
Run the construction at internal accuracy $\rho=c_0\varepsilon$ for a small
constant $c_0$, choose $a=c_1\rho$, and set
\begin{equation}
 M=\log(1/\sin^2a)=O(\log(1/\rho)),\qquad
 \zeta=\Theta(\rho),\qquad
 \gamma=\Theta(\rho^2/M).
 \label{eq:steps}
\end{equation}
By Lemmas~\ref{lem:oracle} and \ref{lem:grid}, the normalized oracle
tolerances are
\begin{equation}
 \epsilon_1=O(\rho),\qquad
 \epsilon_2=O(\rho^2/M).
 \label{eq:tolerances}
\end{equation}
All tested payoffs are bounded by a universal constant after normalization.
Projected gradient ascent over the calibration cube gives average regret
\begin{equation}
 \frac{\operatorname{Reg}_{\rm cal}(T)}T
 =O(\rho)+O\left(\frac{|\mathcal G|}{\rho T}\right).
 \label{eq:cal-regret}
\end{equation}

The imbalanced meta bound has a useful asymmetry
\citep{hu2026optimal}. Since $w_1$ is the loss vertex, the calibration
comparator is at constant squared distance and the loss comparator is at
zero distance. Therefore
\begin{align}
 \mathbb{E}\calK_1
 &\leq O(\rho)
 +O\left(\frac{|\mathcal G|}{\rho T}\right)
 +O(\gamma)+O\left(\frac1{\gamma T}\right),
 \label{eq:upper-cal-bound}\\
 \mathbb{E}\frac1T\sum_t
 [\ell_{\log}(p_t,y_t)-\ell_{\log}(\bar q_t,y_t)]
 &\leq O(\rho^2)+O(M\gamma).
 \label{eq:upper-loss-bound}
\end{align}
The loss line has no $M/(\gamma T)$ term. This is the reason normalization
removes one logarithmic factor rather than merely moving it.
Using $|\mathcal G|=O(\rho^{-1})$ and
$T=\Omega(M\rho^{-3})$ makes
Equations~(\ref{eq:upper-cal-bound}) and
(\ref{eq:upper-loss-bound}) respectively $O(\rho)$ and $O(\rho^2)$.
Equation~(\ref{eq:clip}) adds $O(\rho^2)$ when replacing $\bar q_t$ by
$q_t$. Constant rescaling gives the theorem.
Appendix~\ref{app:upper-proof} includes the complete regret calculation and
the conditional-unbiasedness argument for sampled reports.

\paragraph{Why a direct Lipschitz reduction fails.}
To keep the clipping cost below $\varepsilon^2$, direct clipping needs
$\delta=\Theta(\varepsilon^2)$. Log loss is then
$L=\Theta(\varepsilon^{-2})$-Lipschitz. The
$O(L^2\varepsilon^{-3})$ horizon of \citet{hu2026optimal} becomes
$O(\varepsilon^{-7})$. In contrast, our oracle depends on adjacent KL
tangent gaps, which remain $O(\varepsilon^2)$ under the boundary-refined
grid.

\paragraph{Computational cost.}
The grid has $O(\varepsilon^{-1})$ points. Evaluating all point and adjacent
segment candidates takes $O(\varepsilon^{-1})$ time per round and the chosen
action has support at most two. Appendix~\ref{app:implementation} describes
the exact minimax scan used in our implementation.

\section{A lower bound for log loss}
\label{sec:lower}

The logarithmic factor in Theorem~\ref{thm:upper} is the only remaining gap.
The $\varepsilon^{-3}$ dependence itself is necessary even far from the
boundary.

\begin{theorem}[Lower bound]
\label{thm:lower}
For all sufficiently small $\varepsilon$, any online learner that guarantees
expected $O(\varepsilon)$ ordinary $\ell_1$ calibration and expected
$O(\varepsilon^2)$ excess Bernoulli log loss on every sequence requires
\begin{equation}
 T=\Omega(\varepsilon^{-3}).
 \label{eq:lower-rate}
\end{equation}
This holds even when every hint lies in $[1/4,3/4]$.
\end{theorem}

\paragraph{Proof.}
Use the hard construction of \citet{hu2026optimal}. The hints cycle through
an $\varepsilon$-net of $[1/4,3/4]$, and conditional on the past,
$y_t\sim\Ber(q_t)$. For every non-anticipating prediction $p_t$,
\begin{align}
 \mathbb{E}_t[
 \ell_{\log}(p_t,y_t)-\ell_{\log}(q_t,y_t)]
 &=\KL(\Ber(q_t)\Vert\Ber(p_t))\nonumber\\
 &\geq2(q_t-p_t)^2.
 \label{eq:pinsker}
\end{align}
Thus expected log excess $O(\varepsilon^2)$ implies the same
mean-squared-deviation premise, up to a constant, that drives their rounding
and miscalibration proof. Round $p_t$ to the nearest point of their
$\varepsilon$-net. Calibration increases by at most $\varepsilon$, while
the average expected squared deviation remains $O(\varepsilon^2)$.
Their bin-fluctuation lower bound then gives
$T=\Omega(\varepsilon^{-3})$. This reduction uses the hard instance and
proof, not merely the statement of their squared-loss theorem.

\paragraph{Interpretation.}
The hard hints are bounded away from zero and one. Boundary singularity
therefore does not create the lower bound. It creates the logarithmic gap in
our upper bound through the clipped range $M$. Removing that factor appears
to require either a range-free meta analysis or a different coupling of the
two objectives.

\section{Experiments}
\label{sec:experiments}

The theory is worst-case and its prescribed horizons have large universal
constants. Our experiments instead test three finite-time questions:
whether \method preserves log loss while improving ordinary calibration,
whether the Fisher grid helps on boundary-sensitive sequences, and whether
the same behavior occurs for predictions from a public frozen
representation under shift.

\paragraph{Metrics.}
We report $\calK_1$ from Equation~(\ref{eq:k1}), 20-bin expected calibration
error (ECE), mean log loss, and excess log loss relative to the original
hints. Unlike ECE, $\calK_1$ evaluates every distinct finite-grid report.
All online methods see each outcome only after predicting it.

\paragraph{Methods.}
We compare the unmodified hints, nearest-grid online-SGD Platt scaling
\citep{platt1999probabilistic}, HOPS \citep{gupta2023online} using the same
Platt expert, a prefix histogram recalibrator, calibration-only
approachability, the same two-objective algorithm on a probability-uniform
grid of equal size, and \method on the Fisher grid. Rounding the Platt and
histogram outputs to the evaluation grid makes their ordinary $\calK_1$
finite and directly comparable. HOPS already has ten discrete outputs. The
calibration-only ablation sets the meta weight on loss to zero. We report
both the conservative theorem-motivated \method step and a fixed
$16\times$ practical step selected on two controlled development streams.

\subsection{Controlled streams}

We use six streams that isolate different failure modes: calibrated hints,
a temperature distortion, a boundary shift, an oscillating distortion, a
severe temperature shift, and hints placed exactly on Fisher-grid nodes.
Each stream has $30{,}000$ rounds and five independent label and learner
seeds. The boundary stream includes probabilities near zero and one. The
node stream tests whether a geometry-aware method can leave already
well-positioned hints unchanged.

\begin{table}[t]
\centering
\caption{Controlled online streams at $\varepsilon=0.1$. Each entry is the
macro average over six streams, each with five seeds and $30{,}000$ rounds.
Lower is better for $\calK_1$. Excess is relative to the original hints.}
\label{tab:synthetic}
\small
\begin{tabular}{lrrr}
\toprule
Method & Mean $\calK_1$ & Mean ECE & Mean log excess\\
\midrule
Hints & .0610 & .0604 & $0$\\
SGD Platt & .0199 & .0136 & -.1327\\
HOPS-SGD & .0162 & .0162 & -.1241\\
Prefix histogram & .0183 & .0147 & -.1183\\
Calibration only & \textbf{.0125} & \textbf{.0103} & .2421\\
Equal-size uniform & .0456 & .0442 & -.0628\\
\method (base) & .0479 & .0465 & -.0578\\
\method ($16\times$) & .0233 & .0218 & -.0975\\
\bottomrule
\end{tabular}
\end{table}

The two-objective methods improve average log loss without the large cost of
the calibration-only ablation. At the conservative theorem-motivated meta
step, however, the base method corrects miscalibration more slowly than the
specialized Platt, HOPS, and histogram baselines. The fixed practical step
cuts its macro $\calK_1$ from $0.0479$ to $0.0233$ while increasing the
log-loss improvement. The Fisher-node diagnostic is especially sharp. The
hints have $\calK_1=0.0038$. Both \method variants retain zero log excess
and the same $\calK_1$ to four digits. The equal-size uniform grid
incurs $0.0010$ excess, while calibration-only approachability incurs
$0.5872$. This is not evidence that a nonuniform grid always has smaller
probability mesh. It verifies the intended boundary-preserving geometry
when the hints already lie on its nodes. Appendix~\ref{app:experiments}
also reports a prespecified sweep of the meta step, which traces the expected
calibration-loss frontier. On the severe-temperature stream, multiplying
the base step by 16 reduces \method's $\calK_1$ from $0.1217$ to $0.0234$
while increasing its log-loss improvement from $0.2278$ to $0.3486$.

\subsection{Shifted CIFAR-10 binary stream}

We form a public cat-versus-dog task from CIFAR-10
\citep{krizhevsky2009cifar}. A logistic head is fit on $\ell_2$-normalized,
frozen CLIP ViT-B/32 features \citep{radford2021clip} from the training
split. It reaches $93.95\%$ accuracy on the $2{,}000$ relevant test images.
We repeatedly present random permutations of those test images. Halfway
through the stream, each label is independently flipped with probability
$0.2$, while the original classifier probability remains the hint. This
creates an abrupt calibration shift without retraining the representation or
hint model.

\begin{table}[t]
\centering
\caption{Shifted CIFAR-10 cat-versus-dog stream at $\varepsilon=0.1$.
The full run uses $200{,}000$ rounds per seed and five seeds. Entries are
mean $\pm$ standard error.}
\label{tab:cifar}
\small
\begin{tabular}{lrrrr}
\toprule
Method & $\calK_1$ & ECE & Log loss & Log excess\\
\midrule
Hints & $.1164\mathbin{\pm}.0001$ & $.0563\mathbin{\pm}.0003$
& $.4498\mathbin{\pm}.0007$ & $0$\\
SGD Platt & $.0326\mathbin{\pm}.0002$ & $.0292\mathbin{\pm}.0003$
& $.3533\mathbin{\pm}.0003$ & $-.0965\mathbin{\pm}.0004$\\
HOPS-SGD & $.0276\mathbin{\pm}.0005$ & $.0276\mathbin{\pm}.0005$
& $.3657\mathbin{\pm}.0003$ & $-.0842\mathbin{\pm}.0004$\\
Prefix histogram & $.0613\mathbin{\pm}.0001$ & $.0611\mathbin{\pm}.0001$
& $.3972\mathbin{\pm}.0004$ & $-.0526\mathbin{\pm}.0003$\\
Calibration only & $\mathbf{.0066}\mathbin{\pm}.0001$
& $\mathbf{.0058}\mathbin{\pm}.0005$
& $.6942\mathbin{\pm}.0000$ & $.2444\mathbin{\pm}.0007$\\
Equal-size uniform & $.0162\mathbin{\pm}.0005$ & $.0146\mathbin{\pm}.0006$
& $.3734\mathbin{\pm}.0002$ & $-.0764\mathbin{\pm}.0006$\\
\method (base) & $.0235\mathbin{\pm}.0003$ & $.0218\mathbin{\pm}.0005$
& $.3890\mathbin{\pm}.0004$ & $-.0608\mathbin{\pm}.0004$\\
\method ($16\times$) & $.0083\mathbin{\pm}.0001$ & $.0079\mathbin{\pm}.0001$
& $.3758\mathbin{\pm}.0005$ & $-.0740\mathbin{\pm}.0001$\\
\bottomrule
\end{tabular}
\end{table}

All two-objective methods improve both calibration and log loss after the
shift. The base \method reduces $\calK_1$ by $79.8\%$. The fixed practical
step reduces it by $92.8\%$ and improves mean log loss by $0.0740$ relative
to the frozen hints. Calibration-only approachability
achieves the smallest calibration error but pays $0.2444$ excess log loss,
which exposes the need for the second objective. SGD Platt scaling is
best in log loss on this repeated, one-dimensional shift, followed by its
HOPS wrapper. The equal-size uniform ablation is stronger than the base
\method but weaker than its fixed practical variant in calibration. Most
hints are interior, where uniform-grid probability cells are smaller. This result
distinguishes the theorem's boundary-uniform guarantee from an empirical
claim that Fisher spacing dominates on every distribution.

All settings, streams, metrics, and seeds were fixed before the full runs.
The $16\times$ variant was selected from the prespecified two-stream sweep
and then held fixed. Appendix~\ref{app:experiments} gives the stream
equations, hyperparameters, per-seed values, and a short-run diagnostic used
only to check the pipeline.

\section{Related work}
\label{sec:related}

\paragraph{Recalibration.}
\citet{kuleshov2017estimating} use one calibration subroutine per hint bin
and preserve bounded proper losses. \citet{okoroafor2024faster} introduce a
faster approachability construction for Lipschitz proper losses.
\citet{hu2026optimal} develop imbalanced simultaneous approachability,
obtain the optimal $\Theta(\varepsilon^{-3})$ horizon for globally
Lipschitz losses, and prove the hard instance used in
Theorem~\ref{thm:lower}. Their theorem excludes log loss, and they leave the
need for Lipschitzness open. Our result changes their oracle geometry and
normalizes its unbounded loss coordinate.

\paragraph{Calibration, swap regret, and calibeating.}
Online calibration originates with \citet{foster1998asymptotic}.
\citet{gupta2023online} wrap online Platt scaling with Foster's randomized
forecaster. Because HOPS reports bin midpoints, its calibration error is
ordinary $\calK_1$ on those outputs. It pairs a rate of order $T^{-1/4}$
after tuning with a sharpness comparison, not with realized log-loss
preservation relative to arbitrary hints.
\citet{foster2023calibeating} study calibration while beating the refinement
of other forecasts. Their long version explicitly treats the log score with
pseudocount regularization and finite log grids, so simultaneous log
calibration and log calibeating are not new in a qualitative sense.
\citet{chen2026calibeating} give a no-regret treatment and an optimal finite
action log-calibeating bound, while their simultaneous construction focuses
on Brier loss. \citet{foster2026proper} provide a general proper-calibeating
theory. \citet{fichtl2026bregman} extend calibeating through
Bregman-divergence regret identities. Their loss-based FTRL handles
unbounded log loss, and multiplicative binning handles arbitrary external
forecasts. This is direct prior art for arbitrary-forecast log calibeating,
but it does not give ordinary $\calK_1$ calibration of the replacement
predictions together with a near-optimal
$(\varepsilon,\varepsilon^2)$ finite-horizon guarantee. These objectives
differ from near-optimal ordinary recalibration with realized loss
preservation.

\citet{luo2025kl} connect KL calibration to log-loss swap regret and use a
Fisher-uniform grid with specialized randomized rounding.
\citet{fishelson2025full} relate full swap regret to several discretized
calibration notions. \citet{liu2026adaptive} use the same angular geometry
for pseudo-KL calibration under stochastic nonstationarity. None of these
methods takes an unrestricted external hint and simultaneously preserves its
realized log loss.

\paragraph{Proper losses and boundary-adapted geometry.}
The derivative and Bregman structure of binary proper losses is developed by
\citet{reid2010composite}, while \citet{vanerven2012mixability} characterize
mixability through curvature relative to log loss. The grid used here goes
back to \citet{derooij2009switching} and was also used for entropic online
isotonic regression by \citet{kotlowski2016isotonic}. Recent work on
U-calibration adapts to losses smooth relative to a log barrier, but covers
bounded proper losses rather than log loss itself
\citep{frongillo2026ucalibration}. Offline proper-score decompositions
interpret log-loss improvement through reliability and retained information
\citep{charpentier2026decomposing}. Our setting is adversarial, online, and
finite horizon.

\section{Limitations and conclusion}
\label{sec:conclusion}

The guarantee is in expectation. Sampling the sparse action is needed for
ordinary calibration. Near the boundary, a rare label can make the
difference between two adjacent realized log losses constant even though
their expected tangent gap is $O(\varepsilon^2)$. A range-only martingale
bound would therefore require roughly $\varepsilon^{-4}$ rounds for a
high-probability loss guarantee. A variance-sensitive argument may improve
this, but we do not claim one.

The remaining logarithmic factor may not be fundamental. The lower bound
uses interior hints and gives $\Omega(\varepsilon^{-3})$. We also restrict
the theorem to binary outcomes. Multiclass information geometry suggests an
extension, but sparse simultaneous rounding near simplex faces is a separate
problem.

Within these boundaries, the result shows that unbounded log loss does not
require a polynomial deterioration in online recalibration. The correct
local quantity is adjacent Bayes regret, not a global derivative bound.
Combining this curvature-adaptive oracle with an asymmetrically initialized
meta learner yields a near-optimal arbitrary-hint guarantee.

\section*{Reproducibility statement}

Appendices~\ref{app:oracle-proof}, \ref{app:algorithm}, and
\ref{app:implementation} specify the oracle and algorithm.
Appendix~\ref{app:experiments} specifies data
provenance, stream construction, baselines, hyperparameters, and seeds. The
supplement includes source code, deterministic oracle tests, and machine
readable raw results.

\bibliography{references}
\bibliographystyle{iclr2027_conference}

\appendix

\section{Proof of the nonuniform oracle}
\label{app:oracle-proof}

We first record the proper-loss geometry used by the oracle. For a
differentiable binary proper loss $\ell$, let
$H(p)=(1-p)\ell(p,0)+p\ell(p,1)$ be its Bayes risk and let
$w(p)=-H''(p)\geq0$ be its weight function. Then
\begin{equation}
 \ell_0'(p)=p\,w(p),\qquad
 \ell_1'(p)=-(1-p)w(p).
 \label{eq:proper-derivatives}
\end{equation}
These identities hold almost everywhere under the usual weak regularity
conditions \citep{reid2010composite}. The log loss has
$w(p)=1/[p(1-p)]$.

\begin{lemma}[Adjacent support gap]
\label{lem:general-gap}
Let $0<p<r<1$. The vertical gap from the loss curve at $r$ to the support
line at $p$, whose direction is $(p,p-1)$, equals
\begin{equation}
 G_{\uparrow}(p,r)
 =\frac1p\int_p^r(t-p)w(t)\,dt.
 \label{eq:general-upper-gap}
\end{equation}
After exchanging the two labels, the corresponding gap equals
\begin{equation}
 G_{\downarrow}(p,r)
 =\frac1{1-r}\int_p^r(r-t)w(t)\,dt.
 \label{eq:general-lower-gap}
\end{equation}
For log loss, these quantities are
\begin{equation}
 G_{\uparrow}(p,r)=\frac{\KL(\Ber(p)\Vert\Ber(r))}{p},
 \qquad
 G_{\downarrow}(p,r)=
 \frac{\KL(\Ber(r)\Vert\Ber(p))}{1-r}.
 \label{eq:log-gaps}
\end{equation}
\end{lemma}

\begin{proof}
The support line at $p$ has slope $(p-1)/p$. Its value at the first
coordinate of the curve at $r$ is
\[
 \ell(p,1)+\frac{p-1}{p}[\ell(r,0)-\ell(p,0)].
\]
Subtracting this from $\ell(r,1)$ and applying
Equation~(\ref{eq:proper-derivatives}) gives
\begin{align*}
 &\ell(r,1)-\ell(p,1)
 -\frac{p-1}{p}[\ell(r,0)-\ell(p,0)]\\
 &\quad=\int_p^r
 \left[-(1-t)+\frac{1-p}{p}t\right]w(t)\,dt
 =\frac1p\int_p^r(t-p)w(t)\,dt.
\end{align*}
Reflecting $p\mapsto1-p$ and exchanging labels proves the second identity.
Substitution of $w(t)=1/[t(1-t)]$ and direct integration give
Equation~(\ref{eq:log-gaps}).
\end{proof}

\begin{proof}[Proof of Lemma~\ref{lem:oracle}]
It is useful to work with the unscaled two-label vectors
\begin{equation}
 F_i=\alpha u_i(p_i,p_i-1)
 +\frac{\beta}{M}
 \begin{pmatrix}
 \ell_{\log}(p_i,0)-\ell_{\log}(\bar q,0)\\
 \ell_{\log}(p_i,1)-\ell_{\log}(\bar q,1)
 \end{pmatrix}.
 \label{eq:Fi}
\end{equation}
The desired action is a convex combination of these vectors.

Let $c_i^\perp=(1-p_i,p_i)$. Since
$\langle c_i^\perp,(p_i,p_i-1)\rangle=0$, propriety gives
\begin{align}
 \langle c_i^\perp,F_i\rangle
 &=\frac{\beta}{M}
 \big[
 (1-p_i)\ell_{\log}(p_i,0)+p_i\ell_{\log}(p_i,1)\nonumber\\
 &\hspace{45mm}
 -(1-p_i)\ell_{\log}(\bar q,0)
 -p_i\ell_{\log}(\bar q,1)
 \big]\leq0.
 \label{eq:no-first-quadrant}
\end{align}
Thus no $F_i$ lies strictly in the positive quadrant.

Write
\[
 B=\alpha h(\mathcal G)+\frac{\beta}{M}\tau(\mathcal G).
\]
If some $F_i$ has both coordinates at most $B$, its point mass proves the
claim. Otherwise, every $F_i$ has one coordinate greater than $B$.
Equation~(\ref{eq:no-first-quadrant}) forces its other coordinate to be
nonpositive, so each vector is either upper-left or lower-right.

At the lower endpoint, monotonicity of the partial loss implies
\begin{equation}
 [F_0]_0\leq\alpha p_0\leq\alpha h(\mathcal G).
 \label{eq:left-endpoint}
\end{equation}
Thus its first coordinate is at most $B$. Since no point mass is sufficient,
its second coordinate exceeds $B$, and $F_0$ is upper-left. The symmetric
argument places $F_N$ in the lower-right quadrant. The finite sequence
therefore has an adjacent pair $F_i,F_{i+1}$ that crosses from upper-left
to lower-right.

Suppose first that $p_i\geq1/2$. The line through the shifted and scaled loss
curve at $p_i$ in direction $(p_i,p_i-1)$ is a supporting line.
Lemma~\ref{lem:general-gap} shows that moving from the next curve point to
this line costs vertically
\begin{equation}
 \frac{\beta}{M}
 \frac{\KL(\Ber(p_i)\Vert\Ber(p_{i+1}))}{p_i}.
 \label{eq:upper-cell-cost}
\end{equation}
Let $\widetilde F_{i+1}$ be obtained in two steps. First replace the
calibration direction at $p_{i+1}$ by the direction at $p_i$. Then move the
loss point vertically onto its support line at $p_i$. Both $F_i$ and
$\widetilde F_{i+1}$ now lie on that line, and coordinatewise
\begin{equation}
 \left|[F_{i+1}-\widetilde F_{i+1}]_j\right|
 \leq b_i
 :=\alpha(p_{i+1}-p_i)+
 \frac{\beta}{M}
 \frac{\KL(\Ber(p_i)\Vert\Ber(p_{i+1}))}{p_i}
 \quad (j=0,1).
 \label{eq:modified-endpoint}
\end{equation}
Since $F_{i+1}$ is lower-right, its first coordinate exceeds $B$.
Because $b_i\leq B$, the first coordinate of
$\widetilde F_{i+1}$ is nonnegative. The support-line propriety inequality
forces its second coordinate to be nonpositive. The segment from the
upper-left vector $F_i$ to $\widetilde F_{i+1}$ therefore contains a point
with first coordinate zero and second coordinate nonpositive. Apply the
same mixing coefficient to $F_i$ and the original $F_{i+1}$. By
Equation~(\ref{eq:modified-endpoint}), both coordinates of the resulting
mixture are at most
\begin{equation}
 \alpha(p_{i+1}-p_i)+
 \frac{\beta}{M}
 \frac{\KL(\Ber(p_i)\Vert\Ber(p_{i+1}))}{p_i}.
 \label{eq:upper-crossing-bound}
\end{equation}

For $p_{i+1}\leq1/2$, exchange the labels and reflect the symmetric grid.
Equation~(\ref{eq:upper-crossing-bound}) becomes
\[
 \alpha(p_{i+1}-p_i)+
 \frac{\beta}{M}
 \frac{\KL(\Ber(p_{i+1})\Vert\Ber(p_i))}
 {1-p_{i+1}}.
\]
For a cell crossing $1/2$, use the larger of the two valid oriented bounds.
The definitions of $h(\mathcal G)$ and $\tau(\mathcal G)$ now prove
Equation~(\ref{eq:oracle-bound}).

When $\alpha=0$, take a sequence $\alpha_k\downarrow0$. Every action
constructed above lies in the finite union of closed adjacent edges of the
simplex. A convergent subsequence therefore has an adjacent-support limit,
and continuity on the clipped grid preserves the two inequalities.
\end{proof}

\section{Grid and clipping details}
\label{app:grid-proof}

\begin{proof}[Proof of Lemma~\ref{lem:grid}]
Let the angular spacing be $s_i=\theta_{i+1}-\theta_i\leq a$.
By construction, the grid has
$1+\lceil(\pi/2-2a)/a\rceil=O(a^{-1})$ points.
The mean value theorem and $\sin2\theta\leq1$ give
$p_{i+1}-p_i\leq s_i\leq a$. Also
$p_0=\sin^2a\leq a^2$, and symmetry gives
$1-p_N=p_0$. This proves the first three claims.

For a lower-half cell, write
$p=\sin^2\theta$, $r=\sin^2\phi$, and $s=\phi-\theta$.
The first angle is $a$ and $s\leq a$, hence $s\leq\theta$.
The Bernoulli chi-squared inequality gives the first line of
Equation~(\ref{eq:grid-proof}). The trigonometric identity
\[
 r-p=\sin(\phi+\theta)\sin s
\]
gives its second line. Since $\phi\leq\pi/4$ and
$\phi\leq2\theta$, elementary bounds on sine and cosine yield
\[
 \frac{\sin(\phi+\theta)}
 {\sin\theta\cos\theta}\leq6,
 \qquad
 \frac1{\cos^2\phi}\leq2.
\]
Together with $\sin s\leq s$, this proves the constant $72$. Upper-half
cells follow by reflection. If a cell straddles $1/2$, then its endpoints
lie in $[\sin^2(\pi/8),\sin^2(3\pi/8)]$ for $a\leq\pi/8$, and the same
chi-squared calculation has a universal constant no larger than the stated
loose bound.
\end{proof}

\begin{proof}[Proof of Equation~(\ref{eq:clip})]
Consider $q<\delta$. If $y=1$, clipping reduces the loss. If $y=0$, its
increase is
\[
 -\log(1-\delta)+\log(1-q)
 \leq-\log(1-\delta).
\]
The case $q>1-\delta$ follows by exchanging labels. Since
$\delta=\sin^2a\leq a^2$ and $a\leq\pi/8$,
$-\log(1-\delta)\leq C a^2$ for a universal constant $C$.
\end{proof}

\section{Algorithm and upper-bound proof}
\label{app:algorithm}

\begin{figure}[h]
\centering
\fbox{\begin{minipage}{0.94\linewidth}
\textbf{Algorithm 1: \method}\\[2pt]
\textbf{Input:} target $\rho$, Fisher grid $\mathcal G$, range $M$, steps
$\zeta,\gamma$.\\
\textbf{Initialize:} $u_1=0\in[-1,1]^{\mathcal G}$ and
$w_1=(0,1)\in\simplex^2$.\\
\textbf{For} $t=1,\ldots,T$:
\begin{enumerate}
 \item Observe $q_t$ and clip it to $\bar q_t\in[p_0,p_N]$.
 \item Apply Lemma~\ref{lem:oracle} to
 $(u_t,w_t,\bar q_t)$, obtaining $a_t$ on one point or two adjacent points.
 \item Sample $p_t\sim a_t$ with fresh randomness and report $p_t$.
 \item Observe $y_t$. Form
 $z_t^{(1)}=\{\mathbf1\{p_t=s\}(s-y_t)\}_{s\in\mathcal G}$ and update
 $u_{t+1}=\Pi_{[-1,1]^{\mathcal G}}(u_t+\zeta z_t^{(1)})$.
 \item Let $\bar z_t^{(1)}=\{a_{t,s}(s-y_t)\}_{s\in\mathcal G}$ and
 \[
 \bar z_t^{(2)}=\frac1M\sum_s a_{t,s}
 [\ell_{\log}(s,y_t)-\ell_{\log}(\bar q_t,y_t)].
 \]
 Form
 \[
 g_t=(\langle u_t,\bar z_t^{(1)}\rangle-\epsilon_1,
 \bar z_t^{(2)}-\epsilon_2)
 \]
 and update $w_{t+1}=\Pi_{\simplex^2}(w_t+\gamma g_t)$.
\end{enumerate}
\end{minipage}}
\caption{Operational form of \method. The mixed meta feedback in step 5 is
conditionally equal in expectation to point-mass feedback.}
\label{alg:logrecal}
\end{figure}

\section{Proof of Theorem~\ref{thm:upper}}
\label{app:upper-proof}

Let $d=|\mathcal G|$. Both approachability pairings have magnitude at most
one:
\begin{equation}
 |\langle u,v^{(1)}(a,y)\rangle|\leq1,\qquad
 |v^{(2)}(a,\bar q,y)|\leq1.
 \label{eq:bounded-pairings}
\end{equation}
For the second inequality, every clipped partial log loss lies in
$[0,M]$. The tolerance-subtracted two-vector used by the meta learner
therefore has Euclidean norm at most a universal constant $B$.

Projected online gradient ascent on the cube, initialized at zero and fed
the sampled point-mass vectors $z_t^{(1)}$, satisfies
\begin{equation}
 \sup_{u\in[-1,1]^d}
 \frac1T\sum_{t=1}^T
 \langle u-u_t,z_t^{(1)}\rangle
 \leq\frac{\zeta}{2}+\frac{d}{2\zeta T}.
 \label{eq:cube-regret}
\end{equation}
Here $\|z_t^{(1)}\|_2\leq1$ and
$\|u-u_1\|_2^2\leq d$.

For the meta update, Euclidean online gradient ascent gives, for either
vertex $e_i$,
\begin{equation}
 \frac1T\sum_{t=1}^T\langle g_t,e_i-w_t\rangle
 \leq\frac{\gamma B^2}{2}
 +\frac{\|e_i-e_2\|_2^2}{2\gamma T}.
 \label{eq:meta-regret}
\end{equation}
Conditional on the history through $q_t$ and on the selected mixed action,
the sampled calibration payoff has mean $\bar z_t^{(1)}$. The loss
coordinate in Algorithm~\ref{alg:logrecal} is already the exact mixed mean.
Thus the displayed $g_t$ is conditionally equal to the expected point-mass
residual. The oracle inequality implies
\begin{equation}
 \mathbb{E}[\langle w_t,g_t\rangle\mid\mathcal F_t]\leq0.
 \label{eq:oracle-meta}
\end{equation}
Combining Equations~(\ref{eq:meta-regret}) and
(\ref{eq:oracle-meta}) with $i=1$, then adding
Equation~(\ref{eq:cube-regret}), yields
\begin{equation}
 \mathbb{E}\calK_1
 \leq\epsilon_1+\frac{\zeta}{2}
 +\frac{d}{2\zeta T}
 +\frac{\gamma B^2}{2}+\frac1{\gamma T}.
 \label{eq:cal-complete}
\end{equation}
For $i=2$, the initial divergence vanishes and the loss test has zero
regret, so
\begin{equation}
 \mathbb{E}\frac1T\sum_t
 \frac{\ell_{\log}(p_t,y_t)-\ell_{\log}(\bar q_t,y_t)}M
 \leq\epsilon_2+\frac{\gamma B^2}{2}.
 \label{eq:loss-complete}
\end{equation}

Choose the grid with $a=c_1\rho$. Lemmas~\ref{lem:oracle} and
\ref{lem:grid} give
\[
 d\leq C_1/\rho,\quad
 \epsilon_1\leq C_2\rho,\quad
 M\epsilon_2\leq C_3\rho^2,\quad
 M\leq C_4\log(1/\rho).
\]
Set $\zeta=c_2\rho$ and $\gamma=c_3\rho^2/M$. If
$T\geq C_5M\rho^{-3}$, Equation~(\ref{eq:cal-complete}) is at most
$C_6\rho$ and Equation~(\ref{eq:loss-complete}), after multiplication by
$M$, is at most $C_7\rho^2$. Clipping adds at most $C_8\rho^2$.
Selecting $\rho$ as a sufficiently small constant multiple of the requested
$\varepsilon$ proves both target inequalities.

\paragraph{Adaptive sequences.}
The only stochastic equality in the proof conditions on a history before
the fresh draw $p_t\sim a_t$. Therefore the hints and labels may adapt to
all previous predictions and outcomes. The current outcome may not depend on
the current sampled report. This is the standard expected online
recalibration timing.

\paragraph{Several hint sequences.}
If $m$ hints arrive each round, the aggregating algorithm for one-mixable
log loss produces a single compressed hint whose excess to the best original
hint is at most $(\log m)/T$ \citep{cesabianchi2006prediction}. Applying
Theorem~\ref{thm:upper} to that hint gives the additive horizon
\begin{equation}
 T=O\left(
 \frac{\log(1/\varepsilon)}{\varepsilon^3}
 +\frac{\log m}{\varepsilon^2}
 \right).
 \label{eq:multi-hint}
\end{equation}

\section{Implementation details and deterministic checks}
\label{app:implementation}

For each grid point, the implementation forms the two label payoffs in
Equation~(\ref{eq:mloo-goal}). It evaluates all point masses and every
adjacent segment. On an adjacent edge, the maximum of two affine functions
is minimized either at an endpoint or where the functions cross. The scan
therefore returns the exact minimax action within the union of adjacent
edges in linear time.

The supplement includes the following deterministic checks:
\begin{enumerate}
 \item exact minimax comparisons for random two-label arrays
 \item direct verification of Equation~(\ref{eq:oracle-bound}) on
 $21{,}225$ random oracle inputs over five target accuracies from $0.2$ to
 $0.0125$
 \item support size, adjacency, simplex projection, and sampled-feedback
 checks
 \item a $20{,}000$-round end-to-end run containing hints near both
 boundaries.
\end{enumerate}
Across the same five accuracies, the measured Fisher-grid tangent gap was at
most $0.064\varepsilon^2$. For a probability-uniform grid with the same
number of points and endpoints, the ratio
$\tau/\varepsilon^2$ grew from approximately $1.0$ at
$\varepsilon=0.2$ to $32.1$ at $\varepsilon=0.0125$. These numerical checks
are regression tests and are not used in the proof.

\section{Experimental details and full results}
\label{app:experiments}

\subsection{Controlled stream definitions}

Every controlled stream cycles through a fixed set of hint levels and draws
$y_t\sim\Ber(r_t)$ independently given its truth probability $r_t$.
Let $\sigma$ be the logistic function. The six settings are:
\begin{enumerate}
 \item \textbf{Calibrated:}
 $q_t$ cycles through $\sigma(\{-6,-5.7,\ldots,6\})$ and $r_t=q_t$.
 \item \textbf{Temperature shift:}
 $q_t$ cycles through 41 equally spaced logits in $[-7,7]$. In the second
 half, $r_t=\sigma(\operatorname{logit}(q_t)/2.2+0.55)$.
 \item \textbf{Boundary shift:}
 $q_t$ cycles through
 $\{10^{-7},10^{-5},10^{-3},10^{-2},0.1,0.5,\allowbreak 0.9,0.99,
 0.999,1-10^{-5},1-10^{-7}\}$.
 The truth changes from
 $\operatorname{clip}(0.03+0.94q_t)$ to
 $\operatorname{clip}(0.09+0.82q_t)$ halfway through.
 \item \textbf{Oscillating:}
 31 hint logits in $[-5,5]$ are transformed in blocks of $2{,}000$ rounds
 using temperatures $(0.65,1.8,0.9,2.6)$ and biases
 $(-0.3,0.5,0.2,-0.55)$.
 \item \textbf{Severe temperature:}
 $q_t$ cycles through 41 logits in $[-8,8]$ and
 $r_t=\sigma(\operatorname{logit}(q_t)/3+1)$.
 \item \textbf{Fisher nodes:}
 the first 20, center, and last 20 Fisher-grid points are used as both hints
 and truth.
\end{enumerate}
Truth probabilities in the boundary stream are clipped to
$[10^{-5},1-10^{-5}]$.

\subsection{Hyperparameters and baselines}

The target is $\varepsilon=0.1$. The Fisher grid has 125 points. The
calibration step is $\varepsilon/4$, and the meta step is
$\varepsilon^2/(8M)$. The equal-size uniform baseline shares the Fisher
endpoints and number of points. Calibration-only uses the Fisher grid and
fixes the meta weights to $(1,0)$. We also report a fixed practical
multiplier of 16. It was selected from the prespecified severe-temperature
and boundary-shift sweep, then held fixed for the remaining full runs.

Our online-SGD Platt baseline maintains a slope and intercept on the clipped
hint logit. It uses learning rate
$0.02/\sqrt{1+t/2000}$, clips gradient coordinates to $[-10,10]$, and
rounds each output to the nearest Fisher node. This is not the online Newton
step implementation analyzed by \citet{gupta2023online}. For a direct
comparison to their calibeating wrapper, HOPS runs an independent Foster
1999 randomized game inside each of ten bins of our continuous SGD-Platt
expert, with the expert-bin midpoint as the preferred output. The prefix
histogram has
$\lceil1/\varepsilon\rceil$ hint bins, a one-pseudocount prior at the bin
center, and the same output rounding. These simple adaptive baselines are
included as empirical references, not as implementations of the strongest
possible post-hoc calibration methods.

\begin{table}[h]
\centering
\caption{Mean ordinary $\calK_1$ by controlled stream. Each entry averages
five seeds.}
\label{tab:synthetic-k1-full}
\scriptsize
\begin{tabular}{lrrrrrrrr}
\toprule
Stream & Hints & Platt & HOPS & Histogram & Cal.-only & Uniform & Base & $16\times$\\
\midrule
Calibrated & .0069 & .0140 & .0186 & .0105 & .0120 & .0088 & .0087 & .0090\\
Temperature shift & .0624 & .0148 & .0119 & .0451 & .0124 & .0599 & .0614 & .0346\\
Boundary shift & .0523 & .0479 & .0167 & .0206 & .0124 & .0502 & .0515 & .0289\\
Oscillating & .0440 & .0173 & .0092 & .0190 & .0145 & .0433 & .0442 & .0371\\
Severe temperature & .1968 & .0200 & .0120 & .0117 & .0111 & .1078 & .1180 & .0266\\
Fisher nodes & .0038 & .0053 & .0285 & .0031 & .0128 & .0037 & .0038 & .0038\\
\bottomrule
\end{tabular}
\end{table}

\begin{table}[h]
\centering
\caption{Mean log excess by controlled stream. Each entry averages five
seeds. Negative values improve on the original hints.}
\label{tab:synthetic-excess-full}
\scriptsize
\begin{tabular}{lrrrrrrrr}
\toprule
Stream & Hints & Platt & HOPS & Histogram & Cal.-only & Uniform & Base & $16\times$\\
\midrule
Calibrated & .0000 & .0007 & .0169 & .0097 & .4348 & .0008 & .0001 & .0001\\
Temperature shift & .0000 & -.0765 & -.0632 & -.0422 & .2820 & -.0062 & -.0040 & -.0397\\
Boundary shift & .0000 & -.2501 & -.2609 & -.2624 & .1260 & -.1217 & -.1180 & -.2097\\
Oscillating & .0000 & -.0546 & -.0460 & -.0189 & .2589 & -.0023 & -.0010 & -.0072\\
Severe temperature & .0000 & -.4164 & -.4113 & -.4052 & -.2362 & -.2484 & -.2242 & -.3286\\
Fisher nodes & .0000 & .0006 & .0197 & .0093 & .5872 & .0010 & .0000 & .0000\\
\bottomrule
\end{tabular}
\end{table}

\begin{table}[h]
\centering
\caption{Prespecified Fisher-grid meta-step frontier. The multiplier is
relative to $\varepsilon^2/(8M)$. Entries are means over five seeds.
Parentheses give standard errors.}
\label{tab:frontier}
\scriptsize
\begin{tabular}{rrrrr}
\toprule
& \multicolumn{2}{c}{Severe temperature}
& \multicolumn{2}{c}{Boundary shift}\\
\cmidrule(lr){2-3}\cmidrule(lr){4-5}
Multiplier & $\calK_1$ & Log excess & $\calK_1$ & Log excess\\
\midrule
.25 & .2030 (.0006) & -.0445 (.0006) & .0533 (.0007) & -.1138 (.0019)\\
1 & .1217 (.0005) & -.2278 (.0022) & .0527 (.0005) & -.1204 (.0021)\\
4 & .0561 (.0003) & -.3198 (.0032) & .0414 (.0010) & -.1785 (.0037)\\
16 & \textbf{.0234} (.0007) & \textbf{-.3486} (.0021)
& .0285 (.0011) & -.2169 (.0030)\\
64 & .0289 (.0016) & -.2559 (.0038)
& \textbf{.0179} (.0006) & -.2150 (.0023)\\
\bottomrule
\end{tabular}
\end{table}

The curve confirms that the conservative default does not reflect an
inherent calibration-loss conflict on these streams. Increasing the meta
step first improves both objectives. At the largest multiplier, the meta
weight saturates near calibration and log loss begins to worsen on the
severe-temperature stream. The equal-size uniform grid traces a similar
curve. Its complete values are included in the machine-readable output.
After selecting 16 from these two streams, we ran that multiplier once on
all six controlled streams and on CIFAR-10. Those results appear in
Tables~\ref{tab:synthetic-k1-full}, \ref{tab:synthetic-excess-full}, and
\ref{tab:cifar-seeds}.

The synthetic data seeds are $10{,}000j+s$ for stream index $j$ and
$s\in\{0,\ldots,4\}$. Learner seeds are disjoint offsets. The CIFAR stream
uses five seeds for both permutation and learner randomness.

\subsection{CIFAR-10 provenance}

The features are fixed outputs of the public
\texttt{openai/clip-vit-base-patch32} encoder. We use CIFAR-10 training and
test labels 3 and 5 for cat and dog. The feature vectors are
$\ell_2$-normalized. A logistic regression head with $C=1$, the L-BFGS
solver, and a 2,000-iteration cap is fit once on 10,000 training examples.
It converges in 23 iterations. The 2,000 test examples are permuted and
replayed for 100 passes per seed. Starting after 100,000 examples, each
observed label is flipped independently with probability $0.2$.

\subsection{Per-seed CIFAR-10 results}

\begin{table}[h]
\centering
\caption{Per-seed $\calK_1$ and log excess on the shifted CIFAR-10 stream.}
\label{tab:cifar-seeds}
\small
\begin{tabular}{llrrrrr}
\toprule
& & Seed 0 & Seed 1 & Seed 2 & Seed 3 & Seed 4\\
\midrule
\multirow{2}{*}{Hints}
& $\calK_1$ & .1166 & .1161 & .1161 & .1164 & .1166\\
& Excess & .0000 & .0000 & .0000 & .0000 & .0000\\
\multirow{2}{*}{SGD Platt}
& $\calK_1$ & .0322 & .0331 & .0329 & .0322 & .0324\\
& Excess & -.0961 & -.0973 & -.0961 & -.0956 & -.0977\\
\multirow{2}{*}{HOPS-SGD}
& $\calK_1$ & .0268 & .0286 & .0281 & .0262 & .0283\\
& Excess & -.0833 & -.0851 & -.0840 & -.0831 & -.0853\\
\multirow{2}{*}{Uniform}
& $\calK_1$ & .0165 & .0145 & .0174 & .0172 & .0157\\
& Excess & -.0755 & -.0778 & -.0757 & -.0753 & -.0778\\
\multirow{2}{*}{\method}
& $\calK_1$ & .0236 & .0228 & .0229 & .0239 & .0245\\
& Excess & -.0598 & -.0614 & -.0609 & -.0601 & -.0619\\
\multirow{2}{*}{\method $16\times$}
& $\calK_1$ & .0079 & .0086 & .0084 & .0082 & .0085\\
& Excess & -.0737 & -.0744 & -.0738 & -.0739 & -.0742\\
\bottomrule
\end{tabular}
\end{table}

\section{High-probability boundary}
\label{app:high-prob}

The mixed oracle inequality is pathwise. Ordinary calibration uses sampled
point-mass reports, whose deviation from the mixed calibration vector can be
handled by standard finite-grid martingale concentration at the expected
horizon. Log loss is subtler. Between the first two Fisher nodes, the
realized loss difference on the rare label can be constant even though its
conditional expectation is $O(\varepsilon^2)$. A bound based only on the
range therefore needs $T=\Omega(\varepsilon^{-4})$ to control average
realized loss at scale $\varepsilon^2$.

Possible routes to a near-optimal high-probability theorem include a
conditional-variance analysis tied to the oracle's mixing coefficient and a
correlated rounding scheme. We have not proved either route. All theorems in
the main paper are explicitly in expectation.

\section{ICLR checklist}

\begin{enumerate}
 \item \textbf{Claims:} The upper and lower bounds state their timing,
 expectation, loss, calibration metric, and parameter ranges. Grid novelty,
 a high-probability rate, multiclass extension, and removal of the logarithm
 are not claimed.
 \item \textbf{Proofs:} Full arguments appear in
 Appendices~\ref{app:oracle-proof}--\ref{app:upper-proof}. The only imported
 lower-bound component is identified in Section~\ref{sec:lower}.
 \item \textbf{Experiments:} Data provenance, stream construction,
 hyperparameters, seeds, metrics, and baseline limitations are reported in
 Appendix~\ref{app:experiments}. Code and raw JSON outputs are included in
 the supplement.
 \item \textbf{Compute:} Both reported suites run on a single node. The
 recalibration algorithms use CPUs. One GPU is used only to create the
 frozen CLIP features.
 \item \textbf{Assets:} CIFAR-10 and CLIP are public assets cited in the
 paper. No personally identifying or sensitive data are used.
\end{enumerate}

\end{document}
